% !PRINTED 185
\weberpage{185}

\begin{center}
§. 59.

\gesperrt{Transformation von Formen $n^{\mathrm{ten}}$ Grades.}
\end{center}

Wenn eine Form $n^{\mathrm{ten}}$ Grades von $m$ Veränderlichen, $F(x)$, durch eine lineare Substitution [\S.\,55, (1)] transformirt wird in $\Phi(x')$, so wird gleichzeitig die lineare Function der Veränderlichen $\xi_1, \xi_2 \ldots \xi_m$

\weq{1}{\overset{i}{\varSigma}\, \xi_i F'(x_i)}

durch dieselbe auf $\xi_i$ angewandte Substitution in

% [sic] Im Druck fehlt in (2) der Ableitungsstrich an \Phi, obwohl (1), die Koeffizientenargumentation und die nachfolgende Identität \Phi' verlangen. Der Druck wird unverändert bewahrt.
\weq{2}{\overset{i}{\varSigma}\, \xi_i' \Phi(x_i')}

transformirt, weil beide Ausdrücke den Co\"efficienten der ersten Potenz von $t$ in der Entwickelung von
\[
F(x_i + t\xi_i) = \Phi(x_i' + t\xi_i')
\]
nach Potenzen von $t$ darstellen (\S.\,16). Desgleichen wird die quadratische Function

\weq{3}{\varSigma\, \xi_i \xi_k F''(x_i x_k)}

in

\weq{4}{\varSigma\, \xi_i' \xi_k' \Phi''(x_i', x_k')}

transformirt.

Wenden wir das erste Resultat auf ein System von $m$ Functionen $F_1(x), F_2(x) \ldots F_m(x)$ von beliebigen Graden an und bezeichnen mit $\varDelta$ die Determinante
\[
\begin{vmatrix}
F_1'(x_1),\ F_1'(x_2) \ldots F_1'(x_m) \\
F_2'(x_1),\ F_2'(x_2) \ldots F_2'(x_m) \\
\cdot\ \cdot\ \cdot\ \cdot\ \cdot\ \cdot\ \cdot\ \cdot\ \cdot\ \cdot \\
F_m'(x_1),\ F_m'(x_2) \ldots F_m'(x_m)
\end{vmatrix}\! ,
\]
während wir mit $\varDelta'$ die entsprechende, aus den Functionen $\Phi_1, \Phi_2 \ldots \Phi_m$ gebildete Determinante bezeichnen, so erhalten wir nach \S.\,55, (4), wenn wir wieder mit $r$ die Substitutionsdeterminante bezeichnen,

\weq{5}{\varDelta' = r\varDelta.}

Die Function $\varDelta$ heisst die \gesperrt{Functional-Determinante} der $m$ Functionen $F_1, F_2 \ldots F_m$.

Desgleichen erhalten wir aus der Transformation (3), (4), wenn wir mit $H$ die Determinante
%CONT>>
